\documentclass[11pt, headheight=30pt]{scrartcl}
\usepackage[white, nologo, hw]{darkWhite}

\title{18.745 Lie Algebra}
\subtitle{Homework 1}
\begin{document}
% \maketitle
\section{Connected}
\begin{prob}
    Let $G$ be a connected Lie group and $\tilde{G}$ be its simply connected covering.
    \begin{enumerate}
        \item Show that every discrete normal subgroup of a
        connected Lie group is central.
        \item Apply (1) to $\ker(\tilde{G} \to G)$ to show that
        for any connected Lie group $G$, $\pi_1(G)$ is commutative.
    \end{enumerate}
\end{prob}

In fact, every discrete normal subgroup of a connected topological group is central.

Fix any $h\in H$. Consider the function $f(g) = ghg^{-1}$;
this is a continuous function $G\to G$, which we know maps into $H$ only.
In particular, for all $h'\in H$, for sufficiently small neighborhoods
$h'\in U$, $U\cap H = \{h'\}$; thus the pre-image of each $h'$ must
be a disjoint open cover of $G$. $G$ is connected, thus exactly one
of these must be $G$ and the rest must be $\emptyset$; we
know that $1\in f^{-1}(h)$, so we are done.

We now apply this to $\ker(\tilde{G}\to G)$. As a subgroup of $\tilde{G}$,
this contains exactly the loops at $1$ in $\tilde{G}$, and is naturally
isomorphic to $\pi_1(G)$. Of course, kernels of homomorphisms are
normal subgroups. This is also discrete, because as a covering space,
the pre-image of a sufficiently small neighborhood of $1\in G$ consists
of a \emph{disjoint} union of neighborhoods of elements of $p^{-1}(1)$.
As a central group, $\pi_1(G)$ is abelian.

\section{Morphisms}
% 2. Let f : G → H be a morphism of connected Lie groups such that df : T1G → T1H is an
% isomorphism. Show that f is a covering map, and ker(f) is a discrete central subgroup.
\begin{prob}
    Let $f: G\to H$ be a morphism of connected Lie groups such that
    $df: T_1G\to T_1H$ is an isomorphism.
    Show that $f$ is a covering map, and $\ker(f)$ is
    a discrete central subgroup.
\end{prob}

We first show that $f$ is a covering map. Observe that the open neighborhoods
of $h\in H$ are precisely of the form $hU$ for open neighborhoods $U$ of $1\in H$,
thus it suffices to show that there exists an open neighborhood of $1$
which is evenly covered.

Fine, let's work in coordinates now; $df$ is an isomorphism so $G$ and $H$
have the same dimension. Pick charts $(U, \phi)$ and $(V,\psi)$
around $1\in G$ and $1\in H$ respectively; we care about
$F = \psi\circ f\circ \phi^{-1}$, which is a map from $\phi(U)\subset \RR^n$
to $\psi(V)\in \RR^n$.

If $df: T_1G\to T_1H$ is an isomorphism, this just means that $F$
has total derivative $dF_0: \RR^n \to \RR^n$ which is an isomorphism.
The usual inverse function theorem from multivariable function
says that there exists an open neighborhood of $0\in \RR^n$
mapped homeomorphically to an open neighborhood of $0\in \RR^n$
by $F$. Composing this with $\psi$ and $\phi$ yields open neighborhoods
$U'\subset U$ and $V'\subset V$ containing $1\in H$ and $1\in G$
such that $f$ maps $U'$ homeomorphically to $V'$. In particular, $1\in G$
is isolated from all other points $g\in f^{-1}(1)$, because no other
point in $U'$ maps to $1$.

Left translation is a diffeomorphism on $G$, so any other element
of $g\in f^{-1}(1)\subset G$ admits a neighborhood $gU'$ mapped
homeomorphically to $V'$. All of these can be made disjoint; a brute
force way it to pick a ball inside $U'$ centered at $1$, then take
that ball and shrink it by a factor of $2$ so no two balls overlap.

Thus, $f$ is a covering map. We also showed that $\ker f$
was discrete. By the previous problem, $\ker(f)$ is
a discrete central subgroup.

\section{Representations}
% Let V = {A ∈ M2(C) | A = tA, Tr (A) = 0}. V is a R-vector space of dimension 3.
% Recall
% SU2 = {g ∈ SL2(C) |
% t
% g · g = Id}
% SO3 = {h ∈ SL3(R) |
% th · h = Id}
% (1) Prove SU2 acts on V by conjugation.
% (2) Let ( , ) : V × V → R be the bilinear form on V defined by (A, B) = Tr (AB) for
% A, B ∈ V . Prove that ( , ) is invariant under the action of SU2 in (1).
% (3) Use (1) to define a homomorphism ϕ : SU2 → SO3.
% (4) Show that ϕ is surjective with ker(ϕ) = {±1}.
% (5) Show that SU2 is the universal covering of SO3 and π1(SO3) = Z/2, SO3 ≃ RP3

\begin{prob}
    Let $V = \{A\in M_2(\CC) | A = A^\dagger, \Tr(A) = 0\}$.
    $V$ is a $\RR$-vector space of dimension $3$.
    Recall
    \[
        \SU_2 = \{g\in \SL_2(\CC) | g^\dagger g = I\}
    \]
    \[
        \SO_3 = \{h\in \SL_3(\RR) | h^\top\cdot h = I\}
    \]
    \begin{enumerate}
        \item Prove $\SU_2$ acts on $V$ by conjugation.
        \item Let $(\cdot, \cdot): V\times V\to \RR$ be the bilinear form
        on $V$ defined by $(A,B) = \Tr(AB)$ for $A,B\in V$.
        Prove that $(\cdot, \cdot)$ is invariant under the action of $\SU_2$ in (1).
        \item Use (1) to define a homomorphism $\phi: \SU_2\to \SO_3$.
        \item Show that $\phi$ is surjective with $\ker(\phi) = \{\pm 1\}$.
        \item Show that $\SU_2$ is the universal covering of $\SO_3$ and
        $\pi_1(\SO_3) = \ZZ/2$, $\SO_3 \cong \RP^3$.
    \end{enumerate} 
\end{prob}

\begin{psec*}
    [Action by conjugation]
    Let $g\in \SU_2$ and $A\in V$. We want to show that $gAg^{-1}\in V$.
    \[
        (gAg^{-1})^\dagger = (g^{-1})^\dagger A^\dagger g^\dagger
        = gAg^{-1}
    \]
    \[
        \Tr(gAg^{-1}) = \Tr(g^{-1}gA) = \Tr(A) = 0
    \]
    Next, we want to show that this is a group action.
    \[
        (gh)A(gh)^{-1} = ghA(h^{-1}g^{-1}) = g(hAh^{-1})g^{-1}
    \]
    Actions by $g$ are also obviously linear, so this
    is a representation.
    Finally, we want to show that this is a regular action,
    but we can just follow
    \[
        (g,A)\to (g^{-1}, gA) \to gAg^{-1}
    \]
    which is a composition of regular maps.
\end{psec*}

\begin{psec*}
    [Invariant bilinear form]
    Let $g\in \SU_2$ and $A,B\in V$.
    \begin{align*}
        (gAg^{-1}, gBg^{-1}) &= \Tr(gAg^{-1}gBg^{-1}) = \Tr(gABg^{-1}) \\
        &= \Tr(g^{-1}gAB) = \Tr(AB) = (A,B)
    \end{align*}
\end{psec*}

\begin{psec*}
    [Homomorphism]
    Well, thanks for giving this to us on a silver platter;
    $V$ is a dimension-3 subspace with inner form from $(2)$.

    Conjugation by $\SU_2$ thus corresponds 
    to a unique orthonormal transformations in $\OO_3$.
    This is clearly a homomorphism (everything varies smoothly).
    
    We now claim that this conjugation action by $\SU_3$
    in fact has determinant $1$, which will demonstrate
    that the image lies in $\SO_3$. First observe that by spectral decomposition,
    any element of $\SU_2$ under a suitable change of basis
    can be written as
    \[
        \begin{pmatrix}
            e^{i\theta} & 0\\
            0 & e^{-i\theta}
        \end{pmatrix}.
    \]

    We now pick an orthonormal basis given by ($1/\sqrt{2}$ times)
    the Pauli matrices
    \[
        \sigma_1 = \begin{pmatrix}
            0 & 1 \\ 1 & 0
        \end{pmatrix}
        \quad
        \sigma_2 = \begin{pmatrix}
            0 & -i \\ i & 0
        \end{pmatrix}
        \quad
        \sigma_3 = \begin{pmatrix}
            1 & 0 \\ 0 & -1
        \end{pmatrix}
    \]
    Under conjugation, these become
    \[
        \sigma_1 \to \sigma_1\cos(2\theta) + \sigma_2 \sin(2\theta),\qquad
        \sigma_2 \to -\sigma_1\sin(2\theta) + \sigma_2 \cos(2\theta),\qquad
        \sigma_3 \to \sigma_3
    \]
    This is clearly a rotation with determinant $1$.

    % For verification purposes, we will explicitly write
    % the homomorphism. In Einstein notation, let ${\sigma_i}^\alpha_\beta$
    % be be basis of $V$ (we will use greek letters to
    % denote spinor indices and latin letters to denote
    % vector indices, a la physics; metrics are trivial thus
    % position is aesthetic). Then,
    % \[
    %     2\phi(g)^i_j = g^\alpha_\beta {\sigma^i}^\beta_\gamma (g^{-1})^\gamma_\delta {\sigma_j}^\delta_\alpha
    % \]
    % (we're just extracting the components by taking the dot product, which in this case is
    % a trace, thus a full contraction).
\end{psec*}

\begin{psec*}
    [Double Cover]
    First off, from the above argument,
    $U\in \ker(\phi)$ iff it has eigenvalues $\{1,1\}$ or $\{-1,-1\}$...
    but that just means $U\in \{1, -1\}$. Fine.

    Note that the kernel is $0$-dimensional, yet the dimensions of $\SU_2$
    and $\SO_3$ are both (real) $3$-dimensional. Thus, the image
    of $\phi$ is a $3$-dimensional submanifold of a connected
    $3$-dimensional manifold, thus it must be all of $\SO_3$.
\end{psec*}
    % Okay, now we're going to try to use the results of problem $2$.
    % A local chart of $1\in \SU_2$ is
    % \[
    %     (\theta_1, \theta_2, \theta_3)\to \begin{pmatrix}
    %         e^{i\theta_3} & 0\\
    %         0 & e^{-i\theta_3}
    %     \end{pmatrix},\qquad
    %     \begin{pmatrix}
    %         \cos(\theta_2) & \sin\theta_2\\
    %         -\sin\theta_2 & \cos\theta_2
    %     \end{pmatrix},\qquad
    %     \begin{pmatrix}
    %         \cos(\theta_1) & i\sin\theta_1\\
    %         i\sin\theta_1 & \cos\theta_1
    %     \end{pmatrix}  
    % \]
    % and a local chart of $1\in \SO_3$ is
    % \[
        
    % \]

    % I do some things manually, for fun.

    % We know that elements of $SO(3)$ are generated by
    % rotations in the $X$ and $Z$ directions (from common sense, alternatively
    % use Euler parameters or something)
    % so we'll just explicitly construct those. Above, we generated
    % rotations in the $Z$ direction. In the $X$ direction,
    % we can use the maps
    % \[
    %     \begin{pmatrix}
    %         \cos(\theta) & i\sin\theta\\
    %         i\sin\theta & \cos\theta
    %     \end{pmatrix}.
    % \]
    % Under conjugation,
    % \[
    %     \sigma_1 \to \sigma_1,\qquad
    %     \sigma_2 \to \sigma_2\cos(2\theta) + \sigma_3\sin(2\theta),\qquad
    %     \sigma_3 \to -\sigma_2\sin(2\theta) + \sigma_3\cos(2\theta)
    % \]
    % as desired. Thus, $\phi$ is surjective. To show $\phi$ is a covering
    % map, it suffices to find a neighborhood $U$ of $1\in SO(3)$
    % evenly covered by $\tilde{U}_i$, because then for any other
    % point $h = \phi(g)$, I can just use $hU$,
    % which is evenly covered by $\phi(g)\tilde{U}_i$.
\begin{psec*}
    [Universal Cover]
    First off,
    \[
        \SU_2\cong \{|x|^2+|y|^2=1 \mid x,y \in \mathbb C\}\cong S^3    
    \]
    which is simply-connected.
    Quotienting out by $(x,y)\sim (-x,-y)$ yields $\RP^3$.

    This is a two-sheeted covering, and $\SU_2$ is simply-connected,
    thus $\abs{\pi_1(\SO_3)} = 2$, thus $\pi_1(\SO_3) = \ZZ/2$.
\end{psec*}

\section{Closed Lie Subgroups}
\begin{prob}
    Let $G$ be a Lie group and $H$ a closed Lie subgroup.
    \begin{walk}
        \item Let $\conj{H}$ be the closure of $H$ in $G$.
        Show that $\conj{H}$ is a subgroup in $G$.
        \item Show that each coset $Hx$, $x\in \conj{H}$ is open and dense in $\conj{H}$.
        \item Show that $\conj{H} = H$.
    \end{walk}
\end{prob}

% \begin{claim*}
%     If $N$ is an embedded submanifold of $M$, then $N$ is \vocab{locally closed}.
%     This means that every $x\in N$ has a neighborhood $U\subset M$ such that
%     $U\cap N$ is closed in $U$.
% \end{claim*}

\begin{walk}
    \item
    % Basically, $G$ is not a metric space, but I want to use the word
    % ``sequential.'' Here's how this works: pick some $h\in \conj{H}$, then pick
    % a chart around $h$ in $G$; locally, this looks like $\RR^n$, so there can't
    % be any $\eps$-ball around $h$; this immediately implies that $h$ is locally
    % a sequential limit of $h_i\in H$ under this chart.

    % Okay fine. Now, pick some $h\in \conj{H}$; it is the local sequential
    % limit of $h_i\in H$. Then, $h^{-1}$ is the local sequential limit of
    % $h_i^{-1}$, by sequential continuity of $\ii:H\to G h\mapsto h^{-1}$.

    % Similarly, for any $g,h\in \conj{H}$, $gh$ is the local sequential limit
    % of $g_ih_i$ (because $m:H\times H\to G (g,h)\to gh$ is continuous in $H$).


    % It
    Suppose $h \in \conj{H}$ has some inverse
    $h^{-1} \not\in \conj{H}$. Then $h^{-1}$ has
    a neighborhood $U$ not intersecting $H$. But since the
    inversion map $\iota: G\to G$ sending $h \mapsto h^{-1}$
    is continuous, the preimage $\iota^{-1}(U)$ must be some
    neighborhood of $h \in \conj H$ which doesn't intersect
    $H$, contradicting the fact that $h \in \conj H$.

    Now consider the multiplication map
    $m: \conj H \times \conj H \to G$; assume that
    there's some $h_1h_2=g \not\in \conj H$.
    Draw a neighborhood $U$ of $g$ not intersecting
    $H$. The pre-image is a neighborhood of
    $(h_1,h_2) \in \conj H \times \conj H$ which doesn't
    intersect $H \times H$, which we can assume is of the form
    $V\times W$, but then $V$ is a neighborhood of $h_1$
    not intersecting $H$, contradicting the fact that
    $h_1 \in \conj H$.

    \item %Since $H$ is an embedded submanifold, for any $h \in H$ we can find a neighborhood $U \subseteq G$ such that $\conj H \cap U \subseteq H$. So any point in $H$ has a neighborhood in $\conj H$ which is open in $\conj H$; this implies $H$ is open in $H$.
    $H$ is dense in $\conj H$ by definition.
    We show $H$ is open in $\conj H$. Fix any $h \in H$;
    we will find a neighborhood $U$ of $h$ in $\conj H$
    such that $U \subseteq H$. 
    
    $H$ is a manifold, so we can take a neighborhood $U$ of $h$
    in $H$ such that $\conj U \subseteq H$ (stare at a local chart
    and pick a closed ball).
    
    $U$ is open in the quotient topology, thus there is some open
    $W\subseteq G$ such that $W\cap H = U$. Then, $W \cap \conj U=U$.
    Thus $W \cap \conj H = U$, so $U$ is open in $\conj{H}$.

    % $W \cap \conj U=U$. But this is clear, since we picked $U$
    % such that $\conj U \in H$.
    
    For any other coset, since left translation is a homeomorphism,
    $xW$ is also an open set, thus $xH = xW\cap \conj{H}$ is open
    in $\conj{H}$.
    %First we show that $H$ is open in $\conj H$; for every $h \in H$ we want to find an open set $U$ in $G$ containing $h$, such that $\conj H \cap U \subseteq H$. 

    \item The cosets are supposed to be disjoint but the dense subsets
    intersect any open set! Thus there's only one coset.
\end{walk}

\newpage

\section{Grassmanian}
\begin{prob}
    Let $G_{n,k}$ be the set of $k$-dimensional subspaces of $\RR^n$;
    this is called the \vocab{Grassmanian}.

    \begin{enumerate}
        \item Show that $G_{n,k}$ is a homogenous space of $O_n(\RR)$
        and can thus be identified with $O_n(\RR)/H$ for $H \leq O_n(\RR)$.
        \item Show that $G_{n,k}$ is a smooth manifold of dimension $k(n-k)$.
    \end{enumerate}
\end{prob}
\begin{enumerate}
    \item The action of $O_n(\RR)$ on $\mathbb R^n$,
    naturally extends to act on $k$-dimensional subspaces of $\RR^n$.
    This action is transitivity because for any $k$-dimensional
    subspace, we can select an orthonormal basis and transform
    it into any other orthonormal basis of a $k$-dimensional
    subspace using the action of $O_n(\RR)$. Consequently,
    when a specific element $W \in G_{n,k}$ is fixed, we can
    establish an identification between $G_{n,k}$ and the orbit
    of $W$ under $O_n(\mathbb R)$. This orbit corresponds to
    $O_n(\mathbb R)/H$, where $H$ represents the stabilizer of $W$.

    \item The Grassmannian is a coset manifold $O_n(\mathbb R)/H$.
    The group $O_n$ has dimension $n(n-1)/2$ (the first basis vector
    has $n$ choices, the second has $n-1$, etc.)
    
    The stabilizer $H$ of some $k$-dimensional
    subspace $W$ must preserve both $W$ and $W^\perp$;
    it is a direct product of elements of $O_k$ and $O_{n-k}$.
    Thus the dimension of $O_n(\mathbb R)/H$ is
    \[
        \frac{n(n-1)}{2}-\frac{k(k-1)}{2}-\frac{(n-k)(n-k-1)}{2} = k(n-k)
    \]
\end{enumerate}

\end{document}